\documentclass[11pt,a4paper]{article}
\usepackage[utf8]{inputenc}
\usepackage[spanish]{babel}
\usepackage{amsmath, amssymb}
\usepackage{graphicx}
\usepackage[margin=2.5cm]{geometry}
\usepackage{fancyhdr}
\usepackage{tcolorbox}

\definecolor{UCBblue}{RGB}{0, 51, 153}
\definecolor{UCBgray}{RGB}{100, 100, 100}
\definecolor{UCBred}{RGB}{153, 0, 0}

\pagestyle{fancy}
\fancyhf{}
\fancyhead[L]{\textcolor{UCBblue}{\textbf{Ingeniería Mecatrónica - UCB Tarija}}}
\fancyhead[R]{\textcolor{UCBgray}{Taller de Nivelación - Módulo 4}}
\fancyfoot[C]{\thepage}
\renewcommand{\headrulewidth}{0.4pt}
\renewcommand{\headrule}{\hbox to\headwidth{\color{UCBblue}\leaders\hrule height \headrulewidth\hfill}}

\begin{document}

\begin{center}
    {\Large \textbf{\textcolor{UCBblue}{Guía de Modelado Analítico: Dinámica Frecuencial}}}\\
    \vspace{0.3cm}
    {\large \textbf{Fase CDIO: Concebir (La Necesidad de Automatizar)}}\\
    \vspace{0.5cm}
    \textbf{Estudiante:} \rule{8cm}{0.4pt} \quad \textbf{Fecha:} \rule{3cm}{0.4pt}
\end{center}

\vspace{0.5cm}

\section*{1. Fundamentación: Reactancia Dinámica}
La impedancia total de un sistema mecatrónico no es una constante; es una función de la frecuencia angular ($\omega$) a la que operan sus fuentes. 
$$ \mathbf{Z}_L(\omega) = j\omega L \quad \text{(Inductor)} \quad \quad \mathbf{Z}_C(\omega) = -j\frac{1}{\omega C} \quad \text{(Capacitor)} $$

\section*{2. Misión de Ingeniería: El Límite Humano}
Se le encomienda diseñar un filtro pasivo en serie con $R = 10\Omega$, $L = 10\text{mH}$ ($0.01\text{H}$) y $C = 100\mu\text{F}$ ($0.0001\text{F}$). Voltaje de entrada: $\mathbf{V} = 10\angle0^\circ$ V. 
La impedancia total es $\mathbf{Z}_{eq} = R + j\left(\omega L - \frac{1}{\omega C}\right)$.

Para comprender el espectro de este filtro, usted decide calcular la magnitud de la corriente ($|\mathbf{I}| = \frac{|\mathbf{V}|}{|\mathbf{Z}_{eq}|}$) manualmente en dos extremos frecuenciales.

\vspace{0.3cm}
\textbf{A. Cálculo en Baja Frecuencia ($\omega = 100$ rad/s):}
\begin{tcolorbox}[colback=white,colframe=UCBgray,height=4cm,valign=top]
\textit{Desarrollo analítico para $\omega=100$: Calcule $\mathbf{Z}_{eq}$ y luego $|\mathbf{I}|$.}
\end{tcolorbox}

\textbf{B. Cálculo en Alta Frecuencia ($\omega = 10000$ rad/s):}
\begin{tcolorbox}[colback=white,colframe=UCBgray,height=4cm,valign=top]
\textit{Desarrollo analítico para $\omega=10000$: Calcule $\mathbf{Z}_{eq}$ y luego $|\mathbf{I}|$.}
\end{tcolorbox}

\section*{3. Reflexión (Transición al Código)}
Ha demorado varios minutos en calcular solo \textit{dos} puntos del espectro. Para encontrar la frecuencia exacta de resonancia (donde la corriente es máxima), necesitaríamos calcular miles de puntos intermedios.
\begin{tcolorbox}[colback=red!5!white,colframe=UCBred,title=\textbf{El Propósito de la Programación Modular}]
En ingeniería, no programamos por lujo, sino por supervivencia computacional. En el \texttt{Notebook\_04}, encapsularemos estas ecuaciones en una función (\texttt{def}) de Python y automatizaremos este cálculo tedioso $10,000$ veces por segundo usando \texttt{NumPy}.
\end{tcolorbox}

\end{document}
